\chapter*{Abstract}
This report presents a scalable urban mobility analytics workflow based on the
New York City Yellow Taxi dataset. The solution combines Apache Spark for distributed
processing, Apache Zeppelin for interactive notebooks, and Python/PySpark for
geographic enrichment and map-based visualization.

The implemented pipeline covers end-to-end data handling: Parquet ingestion,
traceable cleaning rules, feature enrichment, and analytical queries with DataFrames
and Spark SQL. Two Zeppelin notebooks organize the work: the first one
(\texttt{Análisis Taxis Nueva York.zpln}) performs loading, cleaning and all
tabular and visual analyses; the second one
(\texttt{Análisis Taxis Nueva York Mapas.zpln}) reads exported zone-level indicators
and generates interactive choropleth maps using Folium.

The resulting system is reproducible in a local Docker-based environment and suitable
for academic presentation. Although it is not a production cluster deployment, the
architecture and code are designed to support future migration to larger Spark
infrastructures without rewriting the analytical logic.

\vspace{.5cm}

\textbf{Keywords:} Big Data, Apache Spark, Apache Zeppelin, PySpark, Parquet,
urban mobility, Spark SQL, data quality, Folium, exploratory analytics.